%%%%%%%%%%%%%%%%%%%%%%%%%%%%%%
%     Open Source
%%%%%%%%%%%%%%%%%%%%%%%%%%%%%%
% Technical variant only. Replaces Content/Personal_Projects/AI_Agents.tex,
% which carries the subagent library as a single line in the other variants.
% If there are less than N\baselineskip (N lines of text) of vertical space left on the current page, the \section{New Section} will start on the next page.
% This ensures that the section heading and at least N lines of text appear together on the same page, avoiding awkward breaks.
\needspace{6\baselineskip}
\runsubsection{Open Source Contributions}
\descript{ }
\location{2024 \textendash{} Ongoing | Upstream Projects}
18 merged pull requests into projects I do not own, three of them security fixes to \href{https://github.com/unraid/ci-runner-farm}{unraid/ci-runner-farm}: a runner registration token passed in Docker's argv, readable out of /proc by any local user, moved to a mode-0600 environment file; secret redaction that covered only the log-viewing paths, moved to one enforcement point over every response reaching the browser; and numeric config values reaching arithmetic contexts unvalidated, validated at parse time. I also taught \href{https://github.com/oxsecurity/megalinter/pull/8686}{MegaLinter's} custom-flavour generator to produce output that passes MegaLinter, a bug I hit building the flavour below and fixed upstream rather than worked around.
\sectionsep{}

\needspace{6\baselineskip}
\runsubsection{Open Source Projects}
\descript{ }
\location{2022 \textendash{} Ongoing | Personal Projects}
\vspace{0.25em}
\begin{tightemize}
  \item \href{https://github.com/laywill/megalinter-custom-flavor-github-latex-markdown}{A custom MegaLinter flavour}: a 25-linter image in place of the all-in-one, published to ghcr.io as a reusable Action, with a daily job that tracks upstream releases and cuts a matching one. It more than halved lint time on this repository, from roughly four and a half minutes to two, and this CV is built by it.
  \item \href{https://github.com/laywill/llm-pdf-renamer}{llm-pdf-renamer} renames PDFs in bulk from their contents using a local LLM, reading embedded text first and falling back to OCR only when there is none. \href{https://github.com/laywill/transcription-tools}{transcription-tools} turns audio and video into subtitles and plain text, and resumes part-way through libraries too large to redo on failure. Both carry four CI workflows and test across Python 3.10 to 3.14.
  \item \href{https://github.com/laywill/awesome-claude-code-subagents}{203 Claude Code subagents across 24 categories}, reworked from an open-source original: each encodes a repeatable engineering task and pins its own tool permissions and model, with categories mapped to five risk tiers by how much a bad run could break.
  \item Since late 2025 the code in my personal projects is AI-generated. I review all of it and own the architecture.
\end{tightemize}
\sectionsep{}
